% Generated by scripts/assemble.py; corrections belong in source/corrections.json.
\examhead{99}{南京航空航天大学 2025 年数学分析试卷}{98}
% p099-b20; source PDF page 99
\par\Needspace{5\baselineskip}
\blockanchor{p099-b20}
\majorquestion{一}求下列极限和积分.\par
% p099-b21; source PDF page 99
\blockanchor{p099-b21}
(1) $\lim_{x\to 0}\left(\frac{e^x + e^{2x} + e^{3x}}{3}\right)^{\frac{1}{x}}.$\par
% p099-b22; source PDF page 99
\blockanchor{p099-b22}
(2) $\int_0^1\frac{e^x(1 - x)^2}{(1 + x^2)^2}\mathrm{d}x.$\par
% p099-b23; source PDF page 99
\par\Needspace{5\baselineskip}
\blockanchor{p099-b23}
\majorquestion{二}设 $f(x)$ 在 $[a, +\infty)$ 上有定义, 证明: $\lim_{x \to +\infty} f(x) = A$ 的充要条件是对任意包含在 $[a, +\infty)$ 上的单调递增数列 $\{x_n\}$, 如果 $\lim_{n \to \infty} x_n = +\infty$, 就有 $\lim_{n \to \infty} f(x_n) = A$.\par
% p099-b24; source PDF page 99
\par\Needspace{5\baselineskip}
\blockanchor{p099-b24}
\majorquestion{三}解答如下问题:\par
% p100-b00; source PDF page 100
\blockanchor{p100-b00}
(1) 设 $f(x)$ 在 $(a, b)$ 上连续, 且 $\lim_{x \to a^{+}} f(x)$ 与 $\lim_{x \to b^{-}} f(x)$ 均存在且有限, 证明: $f(x)$ 在 $[a, b]$ 上一致连续.\par
% p100-b01; source PDF page 100
\blockanchor{p100-b01}
(2) 证明: $f(x) = \frac{|\sin x|}{x}$ 在 (0,1) 与 (-1,0) 上一致连续, 但在 $0 < |x| < 1$ 上不一致连续.\par
% p100-b02; source PDF page 100
\par\Needspace{5\baselineskip}
\blockanchor{p100-b02}
\majorquestion{四}设 $0 < a < b, f(x)$ 在 $[a, b]$ 上连续，在 $(a, b)$ 上可导.\par
% p100-b03; source PDF page 100
\blockanchor{p100-b03}
(1) 证明: 存在 $\xi \in (a, b)$, 使得 $f(b) - f(a) = \xi f'(\xi) \ln \frac{b}{a}$.\par
% p100-b04; source PDF page 100
\blockanchor{p100-b04}
(2) 证明: $\lim_{n\to \infty}n(\sqrt[n]{x} -1) = \ln x.$\par
% p100-b05; source PDF page 100
\par\Needspace{5\baselineskip}
\par\Needspace{6\baselineskip}
\blockanchor{p100-b05}
\majorquestion{五}设 $f(x)$ 在 $[a, b]$ 上二阶可导, 且 $f''(x) \geq 0$, $f(x) \leq 0$. 证明:\par
\nopagebreak[4]
% p100-b06; source PDF page 100
\blockanchor{p100-b06}
\displayformula{f (x) \geq \frac {2}{b - a} \int_ {a} ^ {b} f (x) \mathrm{d} x, x \in [ a, b ].}
% p100-b07; source PDF page 100
\par\Needspace{5\baselineskip}
\par\Needspace{6\baselineskip}
\blockanchor{p100-b07}
\majorquestion{六}设\par
\nopagebreak[4]
% p100-b08; source PDF page 100
\blockanchor{p100-b08}
\displayformula{f _ {n} (x) = \sum_ {k = 1} ^ {n} \frac {1}{n} \cos \left(x + \frac {k}{n}\right), n = 1, 2, \dots}
% p100-b09; source PDF page 100
\blockanchor{p100-b09}
求 $\{f_n(x)\}$ 的极限函数, 并证明 $\{f_n(x)\}$ 在 $(- \infty, +\infty)$ 上一致收敛.\par
% p100-b10; source PDF page 100
\par\Needspace{5\baselineskip}
\blockanchor{p100-b10}
\majorquestion{七}讨论 $\int_{1}^{+\infty} \frac{\sin x \cos \frac{1}{x}}{x} \, \mathrm{d}x$ 是绝对收敛, 条件收敛还是发散?\par
% p100-b11; source PDF page 100
\par\Needspace{5\baselineskip}
\blockanchor{p100-b11}
\majorquestion{八}点 $P(x_0, y_0, z_0)$ 是球面 $x^2 + y^2 + z^2 = 1$ 上的点 $(x_0, y_0, z_0 > 0)$, 求球面在点 $P$ 的切平面与三个坐标平面围成的四面体体积的最小值.\par
% p100-b12; source PDF page 100
\par\Needspace{5\baselineskip}
\blockanchor{p100-b12}
\majorquestion{九}(可能有误) 求 $f(x) = e^{\sqrt{x}}$ 在 $[0, 2\pi]$ 上的傅里叶级数, 并求级数的收敛函数以及 $\sum_{n=1}^{\infty} \frac{1}{n^2 + 2}$ 的值.\par
% p100-b13; source PDF page 100
\par\Needspace{5\baselineskip}
\blockanchor{p100-b13}
\majorquestion{十}$f(x,y)$ 在 $(0,1)$ 附近存在一阶连续的偏导数，且 $f_{y}(0,1) \neq 0, f(0,1) = 0$. 证明：$f\left(x, \int_{0}^{t} \sin u \, \mathrm{d}u\right) = 0$ 在点 $\left(0, \frac{\pi}{2}\right)$ 附近唯一确定单值函数 $t = \varphi(x)$，并求 $\varphi'(0)$.\par
% p100-b14; source PDF page 100
\par\Needspace{5\baselineskip}
\par\Needspace{6\baselineskip}
\blockanchor{p100-b14}
\majorquestion{十一}椭球面 $S: \frac{x^2}{a^2} + \frac{y^2}{b^2} + \frac{z^2}{c^2} = 1 (z \geq 0), (\cos \alpha, \cos \beta, \cos \gamma)$ 是椭球面上的单位外法向量，求曲面积分\par
\nopagebreak[4]
% p100-b15; source PDF page 100
\blockanchor{p100-b15}
\displayformula{\iint_ {S} z \left(\frac {x}{a ^ {2}} \cos \alpha + \frac {y}{b ^ {2}} \cos \beta + \frac {z}{c ^ {2}} \cos \gamma\right) d S.}
% p100-b16; source PDF page 100
\par\Needspace{5\baselineskip}
\blockanchor{p100-b16}
\majorquestion{十二}求重积分 $\iint_{D} \frac{(\sqrt{x} + \sqrt{y})^4}{x^2} \, \mathrm{d}x \, \mathrm{d}y$, 其中 $D$ 为 $\sqrt{x} + \sqrt{y} = 1$, $\sqrt{x} + \sqrt{y} = 2$ 与 $y = x$ 及 $x$ 轴所围区域.\par
